\documentclass{article}
\begin{document}
\begin{abstract}
本文研究带容量约束的配送优化问题，并给出可复算的求解与检验结果。
\end{abstract}
\newcommand{\keywords}[1]{#1}
\keywords{配送优化；整数规划；路径设计}

\section{问题重述}
在给定需求、容量和时间窗的条件下，确定配送顺序与分配数量。

\section{问题分析}
该问题同时包含离散路径与连续数量决策，需统一目标函数、决策变量和约束口径。

\section{模型假设}
需求在一个决策周期内已知，车辆速度保持不变，服务时间计入时间窗。

\section{符号说明}
$x_{ij}$ 表示是否从节点 $i$ 到节点 $j$，$q_i$ 表示节点 $i$ 的配送量。

\section{模型建立与求解}
以总里程最小为目标函数，建立容量约束、流平衡约束和时间窗约束，采用整数规划算法求解。

\section{结果分析}
求得的方案总里程为 128.4 km，所有车辆载重均满足容量约束。相比基准方案，里程下降 11.6\%，这表明路径合并减少了重复折返；其中第二辆车的载重最接近上界，因此该约束是主要活跃约束。结果还说明，节省主要来自相邻客户的连续服务，而不是减少车辆数量。

\section{模型检验}
将决策变量回代全部约束，最大容量残差为 0，时间窗最大误差为 0.3 min，并与人工构造的小规模基准最优解对照。

\section{稳健性分析}
在需求上下扰动 5\% 的情形下重复求解，方案总里程变化不超过 2.1\%，容量约束仍全部满足。

\section{模型评价与改进}
模型能够同时处理路径与装载量，但尚未考虑随机需求；后续可采用情景规划改进。

\begin{thebibliography}{1}
\bibitem{ref1} 某作者. 配送优化方法研究. 2025.
\end{thebibliography}

\appendix
\section{附录：复算说明}
附录记录输入字段、求解器参数和约束回代清单。
\end{document}
